\documentclass{beamer}
\usetheme{Madrid}
\usecolortheme{dolphin}
\definecolor{charcoal}{RGB}{54,69,79}
\setbeamercolor{structure}{fg=charcoal}
\setbeamercolor{palette primary}{bg=charcoal,fg=white}
\setbeamercolor{palette secondary}{bg=charcoal!90,fg=white}
\setbeamercolor{palette tertiary}{bg=charcoal!80,fg=white}
\setbeamercolor{palette quaternary}{bg=charcoal!70,fg=white}
\setbeamercolor{frametitle}{bg=charcoal,fg=white}
\setbeamercolor{title}{bg=charcoal,fg=white}
\usepackage{hyperref}
\usepackage{graphicx}
\usepackage{booktabs}
\usepackage{listings}
\usepackage{color}
\definecolor{dkgreen}{rgb}{0,0.6,0}
\definecolor{gray}{rgb}{0.5,0.5,0.5}
\definecolor{mauve}{rgb}{0.58,0,0.82}
\lstset{
  language=Java,
  aboveskip=3mm,
  belowskip=3mm,
  showstringspaces=false,
  columns=flexible,
  basicstyle={\small\ttfamily},
  numbers=none,
  keywordstyle=\color{blue},
  commentstyle=\color{dkgreen},
  stringstyle=\color{mauve},
  breaklines=true,
  breakatwhitespace=true,
  tabsize=2
}
\title[Sprint 1 Check-In]{Game Project --- Sprint 1 Check-In}
\subtitle{COP2250: Java Programming}
\author{Kevin Pyatt, Ph.D.}
\institute{State College of Florida \\ Pyatt Labs}
\date{Week 13}
\begin{document}
\begin{frame}
  \titlepage
\end{frame}
\begin{frame}{Sprint 1 Requirements}
\textbf{Every team must show today:}
\vspace{0.3cm}
\begin{tabular}{ll}
\toprule
\textbf{Requirement} & \textbf{Evidence} \\
\midrule
3+ distinct classes & Show the files \\
\texttt{main} produces output & Run it live \\
At least one class manages state & Explain it \\
Basic player input handled & Demonstrate it \\
\bottomrule
\end{tabular}
\vspace{0.3cm}
\textit{Compile and run in front of me. That is the only check that matters.}
\end{frame}
\begin{frame}{Demo Protocol}
\textbf{Each team in order:}
\vspace{0.2cm}
\begin{enumerate}
  \item Screen-share terminal
  \item \texttt{ls *.java} --- show all class files
  \item \texttt{javac *.java} --- must compile clean
  \item \texttt{java Main} --- game launches and produces output
  \item Explain which class manages state and why
\end{enumerate}
\vspace{0.3cm}
\textit{If it does not compile, use the remaining class time. Show me before you leave.}
\end{frame}
\begin{frame}{Apply Exception Handling to Your Game}
\textbf{Every game loop needs input protection. Add it now.}
\begin{tabular}{ll}
\toprule
\textbf{Where} & \textbf{What to catch} \\
\midrule
User enters non-integer & \texttt{NumberFormatException} \\
Array or list access & \texttt{IndexOutOfBoundsException} \\
Any unexpected input & \texttt{Exception} as last catch \\
\bottomrule
\end{tabular}
\vspace{0.3cm}
\textit{Exception handling is a Sprint 3 requirement. Start now while the code is simple.}
\end{frame}
\begin{frame}{Sprint 2 Preview}
\textbf{Due Week 14 Thursday:}
\vspace{0.2cm}
\begin{itemize}
  \item At least one abstract class with 2+ concrete subclasses
  \item At least one method override
  \item At least one composition relationship
  \item Core game loop functional
\end{itemize}
\vspace{0.3cm}
\textit{If your game runs and takes input today, you are on track.}\\
\textit{If it does not compile, you are behind. Fix it this weekend.}
\end{frame}
\end{document}
